\section{Model description}
\label{sec:model}

\subsection{Architecture}
\label{sec:architecture}

\policyengineus{} implements United States federal, state, and local tax and
benefit rules as Python formulas operating over a typed entity hierarchy
(person, tax unit, family, household, and program-specific units such as the
marital unit and the \acs{}/\cps{}-style SPM unit), executed on the
\texttt{policyengine-core} simulation engine (a fork of OpenFisca-Core).
\evidence{/Users/maxghenis/PolicyEngine/CLAUDE.md, "Core Infrastructure"
section: "policyengine-core: Fork of OpenFisca-Core that powers country
models and apps"} Parameters (rates, thresholds, income limits) are stored
as YAML, separated from the variable formulas that consume them, so that a
statutory change updates a parameter value with a dated legislative
reference rather than a code change.
\evidence{/Users/maxghenis/PolicyEngine/CLAUDE.md, "Repository Structure
Patterns" and "Parameter Files (YAML)" sections}

\todo{Confirm and cite: is there a published architecture diagram or
description of policyengine-core's simulation model (entity graph,
parameter resolution, vectorized calculation) beyond CLAUDE.md? Check
policyengine-core's own README/docs before v1 submission.}

\subsection{Rules coverage}
\label{sec:coverage}

\todo{Add a generated table (label it \texttt{tab:coverage} and reference it
here as Table~2 or similar once the label resolves) built from
programs.yaml at manuscript-drafting time, not outline time -- the counts
below are what outline drafting (2026-07-04) found and will drift as the
model develops; regenerate before submission, do not hand-copy them into a
table.} The counts below summarize program coverage as recorded in the
model's own coverage registry,
\texttt{policyengine\_us/programs.yaml} (the single source of truth served
via the \texttt{/us/metadata} API and consumed by the public model coverage
page). At \policyengineus{} package version 1.690.7 (the version current
during outline drafting; the version pinned for v1 submission is recorded in
Section~\ref{sec:disclosures}):

\begin{itemize}
\item 67 top-level programs are registered, of which 58 are \texttt{status:
complete}, 6 \texttt{partial}, and 3 \texttt{in\_progress}.
\evidence{/Users/maxghenis/PolicyEngine/policyengine-us/policyengine\_us/programs.yaml,
parsed with PyYAML}
\item By category: Benefits (33), Energy (13), Healthcare (10), Taxes (7),
Education (2), Housing (1), Legal (1).
\evidence{same file, \texttt{category} field}
\item Two entries carry an explicit \taxsim{} validation note in their
metadata: \texttt{federal\_income\_tax} (``Validated against NBER TAXSIM for
2022'') and \texttt{state\_income\_tax} (``Validated against NBER TAXSIM for
2021+'').
\evidence{same file, \texttt{notes} field on those two entries}
\item 4{,}527 parameter YAML files and 4{,}701 variable Python files (excluding
tests and \texttt{\_\_pycache\_\_}).
\evidence{\texttt{find policyengine\_us/parameters -name "*.yaml" | wc -l};
\texttt{find policyengine\_us/variables -name "*.py" | grep -v \_\_pycache\_\_
| grep -v test | wc -l}, run 2026-07-04}
\item State income tax is implemented for 44 states (of the 41 states plus
DC that levy a broad-based individual income tax; \todo{reconcile 44 vs. the
count of states with a statutory income tax -- some entries may be
non-taxing states with a placeholder tree, or the count may include DC and
a small number of special cases; verify against Tax Foundation or state
revenue department listings before asserting a coverage percentage}).
\evidence{\texttt{ls -d policyengine\_us/parameters/gov/states/*/tax/income |
wc -l}, run 2026-07-04}
\end{itemize}

\subsection{Parameter provenance}
\label{sec:provenance}

\todo{Evidence needed. Confirm whether parameter YAML files uniformly carry
a \texttt{reference} field citing the statute or regulation, and if so,
report what fraction do (a spot check or an automated coverage count over
policyengine\_us/parameters). CLAUDE.md states this is the intended
convention (\evidence{"Include legislative references" under Parameter
Files (YAML)}) but the outline-drafting survey did not verify compliance
at scale. This is exactly the kind of load-bearing quantitative claim the
paper must not assert without a generated count.}

\subsection{Versioned coverage over time}
\label{sec:coverage-history}

\todo{Evidence needed. The CHANGELOG (policyengine-us/CHANGELOG.md) records
patch-level releases at roughly weekly cadence as of mid-2026 (version
1.690.7 during outline drafting) and 22{,}437 commits total
\evidence{\texttt{git log --oneline | wc -l}, run 2026-07-04}, but no
existing artifact tracks program-coverage growth over calendar time (e.g.
programs registered per year). If this paper wants a coverage-growth figure,
it must be generated new from git history + programs.yaml snapshots at past
tags -- flag as new analysis, not assembly, if included in v1.}
